\chapter{Methodology}
\label{chap:methodology}

Building upon the theoretical foundations established in Chapter \ref{chap:fundamentals}, this chapter develops the security testing methodology used in this thesis. The methodology translates the cybersecurity obligations of Regulation (EU) 2023/1230 into a traceable and repeatable engineering procedure. It targets the physical IT/OT interface of a safety-related industrial gateway.

The chapter is structured along the derivation chain that forms the scientific core of the methodology. Section \ref{sec:research_design} defines the research design, the constraints, and the scope. Section \ref{subsec:derivation_mapping} derives the fourteen security requirements from the regulation. Section \ref{sec:asset_threat_function_mapping} maps each requirement to the affected assets, the relevant threats, and the security functions that treat them. Section \ref{subsec:traceability_model} consolidates these links into a single traceability chain. Section \ref{sec:approach_selection} selects the verification techniques and assesses the applicability of the company framework. Section \ref{sec:test_derivation} defines the development of the individual test cases. Section \ref{sec:test_environment} defines the test environment and the three network access variants. Section \ref{sec:evaluation_metrics} defines the verification and assessment model that turns test outcomes into requirement verdicts.


\section{Research Design and Methodological Scope}
\label{sec:research_design}

This section defines the research method, the research questions, and the constraints that bound the methodology.

\subsection{Design Science Research Methodology}
\label{subsec:dsrm}

This research follows the Design Science Research Methodology (DSRM) proposed by \citeauthor{Peffers} \cite[p. 46]{Peffers}. DSRM structures the creation and evaluation of an artifact that solves a practical engineering problem. The artifact of this thesis is a standards-based methodology for the functional security testing of an industrial gateway at the IT/OT interface. The six activities of the DSRM process map onto the structure of this thesis:

\begin{enumerate}
	\item \textbf{Problem Identification and Motivation:} The absence of a systematic approach for verifying security functions under Regulation (EU) 2023/1230 is established in Chapter \ref{chap:introduction}.
	\item \textbf{Definition of Objectives:} The objective of a traceable, standards-based verification methodology that uses prEN 50742 and the IEC 62443 series is defined in Chapter \ref{chap:introduction}.
	\item \textbf{Design and Development:} The methodology, its environment, and its governance are designed in this chapter.
	\item \textbf{Demonstration:} The methodology is applied to a physical System Under Test in Chapter \ref{chap:implementation_results}.
	\item \textbf{Evaluation:} The results and the residual validation gap are assessed in Chapter \ref{chap:implementation_results}.
	\item \textbf{Communication:} The results are documented and discussed in this thesis.
\end{enumerate}

\subsection{Research Questions}
\label{subsec:research_questions}

The design and the evaluation of the artifact are guided by three research questions:
\begin{itemize}
	\item \textbf{RQ1:} How can the cybersecurity obligations of Regulation (EU) 2023/1230 be translated into verifiable, component-level security tests for a safety-related industrial gateway?
	\item \textbf{RQ2:} Which tests of the existing company framework are applicable to an embedded gateway, and how can non-applicable tests be identified before execution?
	\item \textbf{RQ3:} What can direct-access testing at the physical IT/OT interface verify about the effectiveness of the implemented security functions, and where are its boundaries?
\end{itemize}


\subsection{Initial Situation, Constraints, and Scope}
\label{subsec:initial_situation}

The System Under Test (SUT) is a safety remote input and output device. It exposes two interfaces with different characteristics. The first interface carries deterministic safety communication over PROFINET and PROFIsafe according to the black-channel principle. The second interface is a diagnostic service reachable over HTTP. The device therefore combines an OT fieldbus path and an IT diagnostic path in a single component.

Ifm operates an established functional test process. This process does not yet verify security functions in a systematic and requirement-driven manner. The methodology developed here addresses this gap by adding a systematic, requirement-driven security verification. It complements the existing functional and safety validation rather than replacing it.

Several constraints shape the methodology and are treated as fixed influencing factors:

\begin{itemize}
	\item \textbf{Black-box target:} The SUT is evaluated as a compiled product. Its source code is not available, which excludes static source analysis.
	\item \textbf{Isolated laboratory:} Testing is performed in a dedicated laboratory. Production systems are excluded.
	\item \textbf{Safety priority:} The safety function must not be degraded by testing. A cybersecurity measure or a test must never introduce impermissible latency into the safety communication \cite[p. 14]{prEN50742_2025}.
	\item \textbf{Draft normative basis:} prEN 50742 is a draft standard. Its clauses are used as technical guidance, not as a certified conformity baseline.
	\item \textbf{Limited practical time:} The practical test phase is time-limited. This affects the number of executed tests.
\end{itemize}

\label{subsec:methodological_scope}
The scope follows from these constraints. The methodology focuses on functional security testing, as distinguished from vulnerability testing in Section \ref{sec:security_testing_approaches}. It verifies whether planned security functions, such as access restrictions or integrity checks, operate as specified \cite[p. 82]{1341418}. Requirement-driven robustness and availability testing is included where a requirement demands it \cite[p. 69]{prEN50742_2025}. This bounded robustness testing is derived from a requirement and differs from open-ended fuzzing that only searches for unknown defects.

Three activities lie outside the scope. Open-ended vulnerability discovery, such as unstructured penetration testing whose goal is to find unknown defects, is excluded. It cannot systematically demonstrate the presence and correct operation of a mandated security function. Static source-code analysis is excluded because no source code is available. Invasive hardware analysis, such as chip-level or side-channel attacks, is excluded because it exceeds a functional verification at the IT/OT interface. The methodology is developed and demonstrated on a single representative SUT.


\section{Requirement Derivation}
\label{subsec:derivation_mapping}
The derivation follows three steps. First, cybersecurity requirements are extracted from Regulation (EU) 2023/1230. This step isolates specific obligations from general legal text and yields the fourteen requirements documented in Appendix \ref{app:sec:mvo_requirements}. Second, each requirement is mapped to the harmonised standard draft prEN 50742. The regulation defines safety goals, whereas prEN 50742 provides technical guidance for machine design. The resulting mapping is documented in Appendix \ref{app:sec:mapping_pren50742}. Third, each requirement is translated into Component Requirements of IEC 62443-4-2. This step produces granular, component-level specifications for industrial network devices, as documented in Appendix \ref{app:sec:mapping_iec62443}.

\section{Asset, Threat and Security Function Mapping}
\label{sec:asset_threat_function_mapping}
This section instantiates the abstract requirement on the concrete device. It maps each requirement to the affected assets, to the threats that target them, and to the security functions that treat those threats.

\subsection{Asset and Threat Mapping}
\label{subsec:asset_threat_mapping}
A standard clause cannot be tested without a concrete target. The methodology therefore adds a mandatory asset-mapping step. The SUT is decomposed into its physical and logical assets, for example hardware interfaces, communication services, and stored configuration data. Each requirement is paired with the assets to which it applies. This step defines not only what must be verified, but also where the corresponding security function is located. The asset mapping can also surface additional component requirements. A requirement may apply to an asset that introduces its own component-level controls, which the direct mapping from the regulation through prEN 50742 to IEC 62443-4-2 does not yield on its own.

Threats are identified with STRIDE, which derives threats from potential attacker goals as described in Section \ref{subsec:stride}. The methodology then examines each security function from two directions. The first direction is the threat that the function mitigates. The second direction is the threat to which the function is itself exposed. A transmission checksum, for example, mitigates accidental corruption, yet it remains exposed to intentional forgery when the checksum is not cryptographic. This distinction prevents a function from being credited without an examination of its own weaknesses.

STRIDE applies at the level of the security requirement rather than to every individual test. A functional baseline or a pure availability demonstration carries no adversarial goal. For such a test case, the target-threat field is set to not applicable and records a functional or availability verification instead of a STRIDE category. A self-inflicted denial-of-service test, in which the tester deliberately exhausts a resource of the device, is the one borderline case and is labelled accordingly.

The final mapping step links each threat to the security function that treats it, which makes the object of the test explicit. A test verifies whether the assigned security function is present and operates correctly for the affected asset, or it documents the gap where the function is absent or only partial.

\subsection{Security Functions}
\label{subsec:security_functions}

A security function is an implemented mechanism that enforces a security objective for one or more assets. The security tests verify such functions, so the relevant categories are defined here once and referenced afterwards. They are grounded in the Foundational Requirements of IEC 62443 and the corresponding component requirements \cite[p. 15]{IEC62443-3-3} \cite[p. 18]{IEC62443-4-2}, and in the technical measures required by prEN 50742 \cite[p. 17]{prEN50742_2025}.
\begin{itemize}
	\item \textbf{Authentication:} verification of the identity of a connecting user, device, or process before access is granted.
	\item \textbf{Authorisation:} restriction of the permitted actions of an authenticated entity, typically through role-based access control and the principle of least privilege.
	\item \textbf{Communication security:} protection of transmitted data through secure protocols and encryption, so that it cannot be read or altered in transit.
	\item \textbf{Integrity protection:} protection of software, firmware, and configuration data against modification, for example through checksums, cryptographic signatures, or a verified boot process.
	\item \textbf{Logging and auditing:} generation and retention of evidence of security-relevant interventions, with sufficient attribution and time reference to be traceable.
	\item \textbf{Availability and resilience:} preservation of the intended function under load, malformed input, or denial-of-service conditions, together with a defined response to detected faults.
\end{itemize}

Not every gateway implements every function, and a function may be absent, partial, or delegated to an organisational control. The methodology therefore examines, for each requirement, which of these functions the device actually provides, and derives a test that verifies the present function or documents the gap. Table \ref{tab:threat_function} relates the STRIDE threat categories to these security functions and provides the bridge from the threat model to the tests.

\begin{table}[H]
	\caption{Mapping of STRIDE Threat Categories to Security Functions}
	\label{tab:threat_function}
	\centering
	\begin{tabular}{p{0.5\textwidth} p{0.5\textwidth} }
		\toprule
		\textbf{Threat Category} & \textbf{Security Function} \\
		\midrule
		Spoofing & Authentication \\
		Tampering & Integrity protection \\
		Repudiation & Logging and auditing \\
		Information Disclosure & Communication security (encryption) \\
		Denial of Service & Availability and resilience \\
		Elevation of Privilege & Authorisation \\
		\bottomrule
	\end{tabular}
\end{table}

The chain from requirement to asset, threat, and security function is the input to the traceability model of the next section.

\section{Traceability Model}
\label{subsec:traceability_model}

The methodology maintains a traceability chain (cf. \cite[pp. 36]{DIN_EN_ISO_IEC_15408_1_2024}, \cite[p. 82]{1341418} \cite[p. 21]{DIN_EN_IEC_62443_4_1_2018}). The chain links a legal obligation to the evidence that supports its verification:

\begin{center}
	Requirement $\rightarrow$ Asset $\rightarrow$ Threat $\rightarrow$ Security Function $\rightarrow$ Test Objective $\rightarrow$ Test Technique $\rightarrow$ Access Prerequisite $\rightarrow$ Evidence $\rightarrow$ Result
\end{center}

The nodes of the chain are defined as follows:
\begin{itemize}
	\item \textbf{Requirement:} The regulatory obligation (for example RQ-011) and its mapped IEC 62443-4-2 control (for example CR 2.1).
	\item \textbf{Asset:} The concrete gateway component to which the requirement applies.
	\item \textbf{Threat:} The attacker goal that targets the asset, derived via \hyperref[subsec:stride]{STRIDE}.
	\item \textbf{Security Function:} The implemented countermeasure, together with any identified gap.
	\item \textbf{Test Objective:} The property to be demonstrated for the security function.
	\item \textbf{Test Technique:} The verification type, classified by the IEC 62443-4-1 SVV categories introduced in Section \ref{subsec:svv}.
	\item \textbf{Access Prerequisite:} The network position required to run the test, expressed as a network access variant (Section \ref{sec:test_environment}).
	\item \textbf{Evidence:} The recorded artifact, such as a scan output, a packet capture, a log export, or a screenshot.
	\item \textbf{Result:} The observed outcome and its assessment against the acceptance criterion.
\end{itemize}

This chain ensures that every executed test serves a distinct requirement and that every requirement can be traced to concrete evidence.

\section{Selection of Verification Techniques}
\label{sec:approach_selection}
The requirements define what must be verified. This section derives how the verification is performed. It evaluates established approaches, assesses the company testing framework, and justifies the selected combination.

\subsection{Evaluation of Approaches}
\label{subsec:approach_evaluation}
Testing approaches are analysed against six criteria: traceability to requirements, availability of concrete test cases, reproducibility, suitability for an embedded OT target, compatibility with functional safety, and evidentiary value. Table \ref{tab:approach_comparison} summarises their main strength and their main limitation.

\begin{table}[H]
	\caption{Comparison of established security testing approaches}
	\label{tab:approach_comparison}
	\centering
	\begin{tabular}{p{0.3\textwidth} p{0.35\textwidth} p{0.35\textwidth}}
		\toprule
		\textbf{Approach} & \textbf{Main strength} & \textbf{Main limitation} \\
		\midrule
		Functional black-box testing \cite[p. 82]{1341418} & Reproducibility and requirement traceability & Does not assess resistance against misuse \\ \addlinespace
		Abstract verification against IEC 62443 \cite[p. 37]{DIN_EN_IEC_62443_4_1_2018} & Defines what must hold at component level & Provides no concrete, executable test cases \\ \addlinespace
		Penetration testing \cite[p. 296]{tauqeer2021analysis} & Realistic attack evidence & Tester-dependent and difficult to reproduce \\ \addlinespace
		Fuzzing and robustness testing \cite[p. 29]{richier2011security} & Suited to error handling and availability & Requires careful monitoring and result interpretation \\ \addlinespace
		Web and IoT testing guides \cite{manana2026pentest} & Procedural depth for HTTP interfaces & Not tailored to OT fieldbus communication \\ \addlinespace
		Risk-based security testing \cite[p. 37]{doCarmo2025} & Prioritises tests by risk & Requires an extensive prior risk analysis \\
		\bottomrule
	\end{tabular}
\end{table}

No single approach satisfies all criteria. Functional testing and abstract verification provide traceability but no attack perspective. Penetration testing and fuzzing provide an attack perspective but weaker reproducibility. The methodology therefore combines a requirement-driven functional core with bounded robustness testing, and it anchors every test type in a recognised verification taxonomy.

\subsection{Applicability Assessment}
\label{subsec:company_framework}
Ifm provides a security testing framework of ten countermeasure-oriented instructions, denoted CM-001 to CM-010. Each instruction binds one countermeasure to one tool from Kali Linux, a Linux distribution that bundles preconfigured security testing tools, and to one acceptance rule. The framework is organised around tools and countermeasures rather than around derived requirements.

The framework cannot be adopted without adaptation. Several instructions assume interfaces that a general IT target exposes but the embedded SUT does not. CM-003 assumes an authentication endpoint, CM-005 assumes a TLS endpoint, and CM-006 assumes an operating-system shell. The SUT provides none of these. This observation motivates an explicit tool-to-target applicability assessment as a step of the methodology. For each planned tool, the methodology checks whether the target exposes the function that the tool assumes.

The applicability assessment produces one of three outcomes. A tool is applicable when the target exposes the assumed interface. A tool is not applicable when the assumed interface is absent, which is an applicability finding and not a product failure. A tool is blocked when an environmental constraint prevents its operation. An absent interface is a property of the device, whereas a blocked tool is a property of the environment. These three outcomes are the minimal distinction that keeps a non-execution from being misread as a failure, because they separate the positive case from the two independent causes of a non-execution.

\subsection{Selected Verification Strategy}
\label{subsec:selected_approach}
The selected methodology combines three elements. The first element is the verification taxonomy of IEC 62443-4-1 Practice SVV, introduced in Section \ref{subsec:svv}, which distinguishes requirements-based testing, threat-mitigation testing, vulnerability testing, and penetration testing \cite[Section 9]{DIN_EN_IEC_62443_4_1_2018}. The taxonomy anchors every test to a defined verification purpose. The second element is a set of requirement-driven test cases, derived from the fourteen requirements and their asset mapping. The third element is a controlled reuse of the company tools, filtered by the applicability assessment.

These three elements correspond to three complementary layers of the derivation. Regulation (EU) 2023/1230 defines why a security property is required and yields the fourteen requirements. IEC 62443-4-2 defines what must hold at the component level through its Component Requirements. IEC 62443-4-1 Practice SVV defines how each property is verified.

The combined approach meets the six criteria of Section \ref{subsec:approach_evaluation}. It preserves requirement traceability, and it yields concrete and reproducible test cases. It retains an attack perspective through the threat-mitigation and the vulnerability tests. It respects the safety priority by bounding the robustness tests to requirement-driven cases. It also remains compatible with the existing industrial validation process, because it reuses the company tools where they apply and documents where they do not.

\section{Test Case Development}
\label{sec:test_derivation}
This section defines how a requirement becomes a documented and repeatable test case. It describes the development process, the specification format, the categorisation, and the evidence model.

\subsection{Test Development Process}
\label{subsec:test_development}
Each test case is derived in three steps. The first step is a critical review of the requirement. The implemented security functions and the identified gaps are cross-checked against the device documentation, so that a test verifies a real mechanism rather than an assumed one. The second step is test planning. For each requirement, one or more scenarios are defined. Each scenario records the property it proves, its technical justification, and its SVV category. The justification refers to established testing guidance, for example NIST SP 800-115 for network-level tests, MITRE ATT\&CK for ICS for adversary techniques, and the OWASP Web Security Testing Guide for the HTTP interface \cite[p. 1-1]{nist_sp800_115} \cite{mitre_attack_ics} \cite{manana2026pentest}. The third step produces the documented test case in a two-level format.

\subsection{Test Specification}
\label{subsec:test_specification}
Every test case is specified on two levels. This separates the product-neutral definition from the environment-specific execution and supports reuse across device variants \cite[p. 37]{Pilorget2012}.

The first level is the abstract test scenario. It is independent of the concrete laboratory and contains the following fields:
\begin{itemize}
	\item \textbf{Traceability:} the requirement, its mapped Component Requirement, and the target asset.
	\item \textbf{Target Threat:} the STRIDE threat that the test exercises. A functional baseline that carries no adversarial goal is marked as not applicable.
	\item \textbf{SVV Category:} the verification type according to IEC 62443-4-1.
	\item \textbf{Objective:} the property to be demonstrated.
	\item \textbf{Steps:} the abstract procedure.
	\item \textbf{Acceptance Criterion:} the exact expected behaviour that constitutes a pass.
\end{itemize}

The second level is the concrete execution. It binds the abstract scenario to the laboratory. It records the required network access variant, the applied tools, and the concrete commands, and it provides the fields for the result and the supporting evidence. The procedural depth of the concrete steps follows NIST SP 800-115 for network-level tests, the OWASP Web Security Testing Guide for the HTTP interface, and MITRE ATT\&CK for ICS for adversary techniques in the OT domain \cite[p. 1-1]{nist_sp800_115} \cite{manana2026pentest} \cite{mitre_attack_ics}. Only representative abstract scenarios are presented in this chapter. The full two-level catalogue is maintained per requirement as a working artefact, and the executed subset with its supporting evidence is reported by security category in Chapter \ref{chap:implementation_results}.

\subsection{Test Categorisation}
\label{subsec:test_categorisation}
Each test case is classified along four dimensions. The first dimension is the security category. It groups tests by the security property under examination and provides the reader-facing structure of the results in Chapter \ref{chap:implementation_results}. The methodology distinguishes seven security categories: identification and asset discovery, communication security, logging and traceability, availability and resilience, robustness, authorisation and access control, and firmware and configuration integrity. Each category groups the tests that verify one of the security functions defined in Section \ref{subsec:security_functions}, and thereby one of the IEC 62443 Foundational Requirements, as Table \ref{tab:category_function_fr} shows. The second dimension is the SVV category, which states the verification purpose. The third dimension is the network access variant, which states the required position on the communication path (Section \ref{sec:test_environment}). The fourth dimension is the execution and evidence status, which states what was achieved for a given test case.

\begin{table}[H]
	\caption{Security categories, the security function they examine, and the IEC 62443 Foundational Requirement}
	\label{tab:category_function_fr}
	\centering
	\begin{tabularx}{\textwidth}{X X l}
		\toprule
		\textbf{Security category} & \textbf{Security function (Sec. \ref{subsec:security_functions})} & \textbf{FR} \\
		\midrule
		Identification and asset discovery & Identification and inventory & FR 1, FR 7 \\ \addlinespace
		Communication security & Communication security & FR 3, FR 4 \\ \addlinespace
		Logging and traceability & Logging and auditing & FR 2, FR 6 \\ \addlinespace
		Availability and resilience & Availability and resilience & FR 7 \\ \addlinespace
		Robustness & Availability and resilience & FR 3, FR 7 \\ \addlinespace
		Authorisation and access control & Authorisation & FR 1, FR 2 \\ \addlinespace
		Firmware and configuration integrity & Integrity protection & FR 3 \\
		\bottomrule
	\end{tabularx}
\end{table}

Table \ref{tab:security_categories} maps each security category to the requirements it primarily addresses. The mapping is not exclusive, because one test activity can support several requirements, and one requirement can appear in more than one category.

\begin{table}[H]
	\caption{Security categories and the requirements they primarily address}
	\label{tab:security_categories}
	\centering
	\begin{tabularx}{\textwidth}{l X}
		\toprule
		\textbf{Security category} & \textbf{Primary requirements} \\
		\midrule
		Identification and asset discovery & RQ-001, RQ-004, RQ-006, RQ-007, RQ-013 \\ \addlinespace
		Communication security & RQ-001, RQ-002, RQ-014 \\ \addlinespace
		Logging and traceability & RQ-003, RQ-008, RQ-009, RQ-012 \\ \addlinespace
		Availability and resilience & RQ-007, RQ-010 \\ \addlinespace
		Robustness & RQ-010 \\ \addlinespace
		Authorisation and access control & RQ-009, RQ-011, RQ-014 \\ \addlinespace
		Firmware and configuration integrity & RQ-002, RQ-005 \\
		\bottomrule
	\end{tabularx}
\end{table}



\section{Security Test Environment}
\label{sec:test_environment}
Security testing can disturb an industrial network. Active scanning or malicious input can interrupt communication and, in the worst case, affect a safety function. In Operational Technology, availability and safety hold the highest priority. Security tests must therefore never run in a production environment. A dedicated and isolated test environment is a mandatory requirement, in line with NIST SP 800-82 \cite[p. 2]{NIST.SP.800-82r3}.


\label{subsec:hybrid_approach} 
The feasibility of a test depends on the position of the test system on the communication path. The methodology defines three network access variants. Each variant provides a different level of access. Figures \ref{fig:nv1}, \ref{fig:nv2}, and \ref{fig:nv3} illustrate the three variants.

\subsection{NV1: Direct Access}
The test system communicates as a regular participant with the SUT and the controller over the switch. NV1 supports all IP-based tests, including the IoT-Core HTTP services, and any test in which the test system acts as an unexpected communication partner.
\begin{figure}[H]
	\centering
	% --- Globale TikZ-Stile für einheitliches Design ---
	\tikzset{
		node distance=1.5cm and 2cm,
		box/.style={
			draw=black!70, thick, rectangle, rounded corners=1.5mm, 
			minimum width=3.8cm, minimum height=1.4cm, align=center, 
			drop shadow={opacity=0.15}, fill=white, font=\footnotesize
		},
		switch/.style={
			box, fill=gray!10, minimum height=1.4cm, minimum width=3.5cm
		},
		link/.style={draw=black!80, thick, <->, >=Latex},
		cyclic/.style={draw=blue!70, thick, <->, >=Latex},
		mirror/.style={draw=red!80, thick, ->, >=Latex, dashed}, 
		inline/.style={draw=orange!90!black, thick, <->, >=Latex} 
	}
	
	\begin{tikzpicture}
		\node[switch] (sw) {\textbf{Industrial Switch}};
		\node[box, left=of sw] (plc) {\textbf{PLC (F-Host)}};
		\node[box, right=of sw] (sut) {\textbf{SUT}};
		\node[box, below=1cm of sw] (kali) {\textbf{Kali Linux}\\(Regular Participant)};
		
		\draw[link] (plc) -- (sw);
		\draw[link] (sut) -- (sw);
		\draw[link] (kali) -- (sw);
		
		\draw[->, >=Latex, dashed, black!60] (kali.north west) to[bend left=15] 
		node[left, font=\scriptsize, align=right] {IP-based tests\\(e.g., HTTP)} (plc.south east);
		\draw[->, >=Latex, dashed, black!60] (kali.north east) to[bend right=15] 
		node[right, font=\scriptsize, align=left] {Unexpected\\comm. partner} (sut.south west);
	\end{tikzpicture}
	\caption{\textbf{NV1, direct access:} Test system acts as a regular network participant.}
	\label{fig:nv1}
\end{figure}

\subsection{NV2: Passive Observation}
A switch mirror or a network tap makes traffic visible that is not addressed to the test system, such as the cyclic PROFINET and PROFIsafe communication between the controller and the SUT. NV2 requires a switch that supports and is configured for port mirroring. A network tap is an alternative to a mirror port. It duplicates the traffic on a dedicated hardware path and avoids the frame loss that a mirror port can introduce under high load. Address resolution spoofing is not a substitute, because the cyclic traffic uses Ethernet-level addressing (EtherType 0x8892 and DCP) rather than IP and ARP.
\begin{figure}[H]
	\centering
	% --- Globale TikZ-Stile für einheitliches Design ---
	\tikzset{
		node distance=1.5cm and 2cm,
		box/.style={
			draw=black!70, thick, rectangle, rounded corners=1.5mm, 
			minimum width=3.8cm, minimum height=1.4cm, align=center, 
			drop shadow={opacity=0.15}, fill=white, font=\footnotesize
		},
		switch/.style={
			box, fill=gray!10, minimum height=1.4cm, minimum width=3.5cm
		},
		link/.style={draw=black!80, thick, <->, >=Latex},
		cyclic/.style={draw=blue!70, thick, <->, >=Latex},
		mirror/.style={draw=red!80, thick, ->, >=Latex, dashed}, 
		inline/.style={draw=orange!90!black, thick, <->, >=Latex} 
	}
	
	\begin{tikzpicture}
		\node[switch] (sw) {\textbf{Industrial Switch}\\(Mirror Port Configured)};
		\node[box, left=of sw] (plc) {\textbf{PLC (F-Host)}};
		\node[box, right=of sw] (sut) {\textbf{SUT}};
		\node[box, below=1cm of sw] (kali) {\textbf{Kali Linux}\\(Passive Listener)};
		
		\draw[cyclic] (plc) -- (sw);
		\draw[cyclic] (sw) -- (sut);
		\draw[link] (kali) -- (sw);
		
		\node[above=0.2cm of sw, font=\scriptsize, text=blue!80] {Cyclic PROFINET/PROFIsafe Traffic};
		
		\draw[mirror] ($(sw.center) + (0.5,0)$) to[bend left=20] 
		node[right, font=\scriptsize] {Mirrored Traffic (L2)} (kali.north);
	\end{tikzpicture}
	\caption{\textbf{NV2, passive observation:} Switch mirrors cyclic traffic to the test system.}
	\label{fig:nv2}
\end{figure}
\subsection{NV3: Active In-Path Manipulation}
The test system is inserted inline between the SUT and the switch and operates as a transparent bridge. NV3 enables the controlled suppression, delay, or replacement of individual frames in transit.

\begin{figure}[H]
	\centering
	% --- Globale TikZ-Stile für einheitliches Design ---
	\tikzset{
		node distance=1.5cm and 2cm,
		box/.style={
			draw=black!70, thick, rectangle, rounded corners=1.5mm, 
			minimum width=3.8cm, minimum height=1.4cm, align=center, 
			drop shadow={opacity=0.15}, fill=white, font=\footnotesize
		},
		switch/.style={
			box, fill=gray!10, minimum height=1.4cm, minimum width=3.5cm
		},
		link/.style={draw=black!80, thick, <->, >=Latex},
		cyclic/.style={draw=blue!70, thick, <->, >=Latex},
		mirror/.style={draw=red!80, thick, ->, >=Latex, dashed}, 
		inline/.style={draw=orange!90!black, thick, <->, >=Latex} 
	}
	
	\begin{tikzpicture}
		% Knotenpunkte exakt an der gleichen Stelle wie in NV1 und NV2
		\node[switch] (sw) {\textbf{Industrial Switch}};
		\node[box, left=of sw] (plc) {\textbf{PLC (F-Host)}};
		\node[box, right=of sw] (sut) {\textbf{SUT}};
		\node[box, below=1cm of sw] (kali) {\textbf{Kali Linux}\\(Transparent Bridge)};
		
		% Unveränderte Verbindung zum PLC
		\draw[link] (plc) -- (sw);
		
		% Visualisierung der "getrennten" Direktverbindung
		\draw[draw=black!30, thick, dashed] (sw) -- (sut) 
		node[midway, above, font=\scriptsize, text=black!50] {Direct link disconnected};
		
		% Der Traffic wird nun physisch nach unten durch das Kali-System geroutet
		% 1. Kabel vom Switch zum Kali
		\draw[link] (sw) -- (kali) 
		node[midway, left, font=\scriptsize, align=right] {Original\\Traffic};
		
		% 2. Kabel vom Kali zum SUT (mit einem sauberen rechten Winkel gezeichnet)
		\draw[inline] (kali.east) -| (sut.south) 
		node[pos=0.25, above, font=\scriptsize, text=orange!90!black, align=center] {Manipulated\\Traffic};
		
		% Erklärungstext unter Kali
		\node[below=0.1cm of kali, font=\scriptsize, text=black!60] {Delay / Drop / Replace frames};
	\end{tikzpicture}
	\caption{\textbf{NV3, active in-path manipulation:} Test system operates inline to manipulate transit frames.}
	\label{fig:nv3}
\end{figure}

%Only NV1 is realised in this thesis. NV2 requires a configured mirror port, and NV3 requires a physical inline insertion. Both remain part of the methodology as design variants and as future validation steps. A test that depends on NV2 or NV3 is therefore specified but not executed. The as-built topology and the reasons for this restriction are documented in Chapter \ref{chap:implementation_results}.


\subsection{Test System and Tooling}
\label{subsec:execution_tooling}
Security testing exercises an object under defined conditions and compares the actual behaviour with the expected behaviour \cite[p. ES-1]{nist_sp800_115}. The SUT is evaluated as a compiled product without source access, only requirement documentation. The methodology therefore uses dynamic black-box testing, also known as Dynamic Application Security Testing. Static source analysis is out of scope, as stated in Section \ref{subsec:methodological_scope} \cite[p. 6777]{aydos_security_2022}.
The methodology adopts Kali Linux as the standardised operating environment. Several options are considered. A general-purpose Linux distribution with individually installed tools offers flexibility, but it increases setup effort and reduces reproducibility. Commercial security suites offer integrated reporting, but their closed nature limits transparency and independent reproduction. Comparable open security distributions, such as Parrot OS, provide a similar toolset. Kali Linux is selected because it bundles a maintained and pre-configured set of open-source auditing tools, which improves reproducibility and removes the need to develop custom exploit code \cite{kali_whatis}. Its use of the same platform as the company framework is a secondary benefit that supports transfer into the industrial process, not the primary reason for the selection. The tools are grouped by function: network discovery and service enumeration, HTTP and API interception, credential testing, and load or robustness testing \cite[p. A-1]{nist_sp800_115}. A tool is used only after the applicability assessment of Section \ref{subsec:company_framework} confirms that the target exposes the assumed interface. 

The engineering tool for the safety configuration, TIA Portal, is used only as a supporting instrument, for example to read identification and maintenance data or to observe the safety qualifier. It is not treated as a security testing tool.










\section{Verification and Assessment Model}
\label{sec:evaluation_metrics}

This section defines how a test outcome becomes a requirement verdict. The model has three parts: the verdict scale for a single test, the aggregation rule for a requirement, and the evidence rule that decides when a test counts at all. The model deliberately avoids a numeric score, because the practical phase executes only a subset of the specified tests, so a numeric score would suggest a completeness that the bounded demonstration does not have. The assessment is therefore qualitative and evidence-based \cite[p. 144, 217]{Witte2019}.

\subsection{Test Verdicts}
\label{subsec:verdicts}

The classification of Section \ref{subsec:test_categorisation} records what a test addresses and what was achieved. The verdict condenses this into a single outcome. Four scored verdicts and two unscored states are distinguished \cite[p. 1553]{11050621} \cite[p. 23]{richier2011security}:
\begin{itemize}
	\item \textbf{PASS:} the test case was executed, all mandatory acceptance criteria are satisfied, and the evidence is sufficient to demonstrate the security function.
	\item \textbf{PARTIAL:} the test case was executed and the mandatory criteria are satisfied, but a non-critical criterion remains unfulfilled or the mechanism showed minor degradation.
	\item \textbf{FAIL:} the test case was executed and a mandatory security function was bypassed, or the device demonstrably deviates from the obligation.
	\item \textbf{NOT APPLICABLE:} the assumed interface or function is absent, or the topology does not support the test. This decision requires a documented technical justification and is not a failure.
	\item \textbf{INCONCLUSIVE:} the test was run but the stored evidence is insufficient for a definitive outcome.
	\item \textbf{BLOCKED:} the test could not be run because of an environmental or tool constraint.
\end{itemize}
Only the four scored verdicts contribute to a requirement assessment. The two unscored states, inconclusive and blocked, record why no scored verdict was reached and keep a non-execution from being misread as a failure.

\subsection{Requirement Assessment}
\label{subsec:requirement_assessment}

The verdicts are aggregated per requirement. A requirement is assessed as FAIL if any associated test case fails or if the device demonstrably deviates from the obligation. It is assessed as PASS only if every associated and applicable test case passes with sufficient evidence. It is assessed as PARTIAL if the executed cases pass or partially pass but the requirement is not yet fully covered, for example because a test that depends on a network access variant was not executed. A requirement for which no associated test produced a scored verdict is reported as not assessed, together with the reason. This scheme is consistent with the partially validated character of the artifact: in this thesis most requirements remain PARTIAL, not assessed, or inconclusive rather than PASS, which reflects the bounded scope of the practical phase rather than a deficiency of the method.

\subsection{Evidence Model}
\label{subsec:evidence_model}

The assessment rests on one evidence rule that is applied without exception: a test is considered executed only if a corresponding stored artefact exists. An artefact is a recorded output of the test, for example a tool log, a packet capture, a log export, or a screenshot. A reported execution without such an artefact is treated as inconclusive, not as a pass. A tool execution is not, by itself, a requirement verification: the observation it produces must be interpreted against the acceptance criterion, and this interpretation is kept separate from the raw observation. Expected behaviour taken from a test specification is never presented as an observed result. This rule is the basis of the evidence-first reporting applied in Chapter \ref{chap:implementation_results}.

The same rule governs the relationship to existing functional and safety tests. The SUT is a safety device, so some security-relevant properties are already exercised by functional or safety validation. Such a test does not, on its own, prove cybersecurity compliance, because safety validation examines behaviour under random faults, whereas a security verdict requires an adversarial perspective and a security-specific acceptance criterion. Safety evidence is therefore reused only when it is re-evaluated against such a criterion. The organisational embedding of these verification activities into a secure development lifecycle, including the governance responsibilities, is described in Appendix \ref{app:SecureDevelopmentLifecycle}.

